\begin{theorem}[structure-sheaf direct images]
\label{thm:structure-direct-images-audit}
\Verified\ For the constructed map \(f:X\to B=\mathbb P^1\),
\[
 f_*\mathcal O_X=\mathcal O_B,\qquad
 R^1f_*\mathcal O_X\simeq
 \mathcal O_B\oplus\mathcal O_B(-1),\qquad
 R^2f_*\mathcal O_X\simeq\mathcal O_B(-1).
\tag{9.12}
\]
All three direct images are locally free and every base-change morphism
\[
 (R^qf_*\mathcal O_X)\otimes k(b)
 \longrightarrow H^q(X_b,\mathcal O_{X_b}),
 \qquad 0\le q\le2,
\tag{9.13}
\]
is an isomorphism, including at the two nonreduced fibres.  Moreover the
invariant class \(\gamma\) defines the trivial summand in
\(R^1f_*\mathcal O_X\) and is nowhere zero on \(B\).
\end{theorem}

\begin{proof}
The proper map has connected fibres and \(B\) is normal, so Stein
factorization gives \(f_*\mathcal O_X=\mathcal O_B\).

The map is flat.  A local parameter on \(B\) pulls back to a nonzero
element of the regular local domain \(\mathcal O_{X,x}\); if it vanished
identically near one point, analytic continuation on the connected
manifold would make \(f\) constant.  Thus \(\mathcal O_{X,x}\) is
torsion-free, hence flat, over the discrete valuation ring
\(\mathcal O_{B,f(x)}\).

Over \(B^\circ\), the fibres are two-dimensional complex tori and their
structure-sheaf cohomology dimensions are \(1,2,1\).  At the cusp the same
dimensions hold by \cref{lem:cusp-structure-sheaf-audit}.  Grauert's
theorem therefore gives local freeness of ranks \(1,2,1\) and base change
on a neighbourhood of \(B^\circ\cup\{p_0\}\); the precise analytic statement is
\cite[Chapter~III, \S4, Theorem~4.12]{BanicaStanasila1976}, originally due to
Grauert \cite{Grauert1960}.

At \(p_j\), use the finite cover \(J'_j\to N_j\).  If
\(\widetilde f:J'_j\to\Delta_j\) is the smooth torus family and
\(\pi(s)=s^{m_j}\), exactness of finite-group invariants in characteristic
zero gives
\[
 R^qf_*\mathcal O_{N_j}
 =\bigl(\pi_*R^q\widetilde f_*\mathcal O_{J'_j}\bigr)^{G_j}.
\tag{9.14}
\]
After averaging, a free eigenframe \(v\) may be chosen for each character
\(\nu\).  The invariant sections are exactly
\(s^rv\,\mathcal O_{D_j}\), where \(0\le r<m_j\) is determined by
\(\nu\zeta_j^r=1\).  Hence (9.14) is free of ranks \(1,2,1\).
Higher direct images vanish by the fibre-dimension bound, so these are all
the nonzero direct images.  Since all direct images are locally free, the analytic
cohomology-and-base-change theorem, applied by descending induction from
\(q=2\), proves (9.13) at the finite fibres as well.  In the notation of
\cite[Chapter~III, \S3, Theorem~3.4, Corollaries~3.5, 3.7, and 3.10]
  {BanicaStanasila1976}.  Corollary~3.10 says that, for each fixed \(q\),
  freeness of \(R^p f_*\mathcal O_X\) for every \(p\ge q\) is equivalent to
  cohomological flatness in all dimensions \(p\ge q\).  Applying this
  successively with \(q=2,1,0\) gives the displayed base-change maps in every
  relevant degree.  This is a local statement on the base and so applies
  equally to the nonreduced scheme-theoretic fibres.

Write \(G=R^1f_*\mathcal O_X\).  We construct its distinguished section
without using a cellular collapse.  On \(B^\circ\), the monodromy-invariant
integral class \(\gamma\) maps to a holomorphic section
\(\gamma^{0,1}\) of \(G\).  It is nowhere zero there: if its image in
\(H^{0,1}(F)\) vanished, then
\(\gamma\in H^{1,0}(F)\cap H^1(F;\mathbb R)=0\).  On a finite cover it is
a constant invariant class and remains nonzero at \(s_j=0\), so it
descends across \(p_j\).

At the cusp, let \(j:\Delta^*\hookrightarrow\Delta\).  The natural
specialization morphism
\[
 R^1f_*\mathbb C\longrightarrow
 j_*R^1f_*\mathbb C|_{\Delta^*}
\]
is an isomorphism: away from \(0\) this is tautological, while at \(0\) it
is exactly the isomorphism from \(H^1(W;\mathbb C)\) onto the
\(T_0\)-invariant subspace proved in
\cref{thm:p0-integral-specialization}.  Hence the invariant classes
\(\gamma,u\) extend as sections of \(R^1f_*\mathbb C\) on the whole disc.
The natural morphism induced by \(\mathbb C\hookrightarrow\mathcal O\),
\[
 R^1f_*\mathbb C\otimes\mathcal O_\Delta
 \longrightarrow G|_\Delta,
\]
sends them at the central point to the two classes in
\(H^1(W,\mathcal O_W)\).  These form a basis by
\cref{lem:cusp-structure-sheaf-audit}.  Thus
\(\gamma^{0,1}\) extends across the cusp and is nowhere zero everywhere.
In particular \(\mathcal O_B\gamma^{0,1}\subset G\) is a line subbundle.

Relative duality applies to the proper flat Gorenstein map \(f\): \(X\) is
smooth, the base is a smooth curve, and \(\omega_{X/B}\) is invertible.
Using \cref{lem:hodge-line-audit}, it gives
\[
 R^2f_*\mathcal O_X\simeq\mathfrak K^\vee
 \simeq\mathcal O_B(-1).
\tag{9.15}
\]
The relative theorem is the one of Ramis--Ruget--Verdier
\cite[pp.~261--283]{RamisRugetVerdier1971}; the absolute singular
dualizing-complex comparison is independently recorded in
\cite[Theorems~1--2 and Proposition~2]{RamisRuget1970}.

The relative cup product
\[
 \bigwedge^2G\longrightarrow R^2f_*\mathcal O_X
\tag{9.16}
\]
is an isomorphism.  It is so on every smooth torus fibre.  At the cusp it
is so by (9.13) and \cref{lem:cusp-structure-sheaf-audit}.  At \(p_j\),
choose on the cover an eigenframe \(v_1,v_2\), with
\(v_1=\gamma^{0,1}\) invariant.  Its invariant exponent is \(r_1=0\).
The action on \(v_1\wedge v_2\) is dual to the Hodge character in (9.3),
so its invariant generator is \(s_j^{a_j}(v_1\wedge v_2)\).
The invariant generator for the second line is consequently
\(s_j^{a_j}v_2\), because \(0<a_j<m_j\).  Wedge product sends
\(v_1\wedge(s_j^{a_j}v_2)\) to exactly the invariant generator of the
right side of (9.16), with no wrap modulo \(m_j\).

Consequently
\(\det G\simeq R^2f_*\mathcal O_X\simeq\mathcal O_B(-1)\).  The quotient
of \(G\) by its trivial subbundle is \(\mathcal O_B(-1)\), and the
extension splits because
\[
 \operatorname{Ext}^1(\mathcal O_B(-1),\mathcal O_B)
 =H^1(B,\mathcal O_B(1))=0.
\]
This proves (9.12) and identifies the trivial summand.
\end{proof}
